\subsection{Entri \texorpdfstring{$(i,j)$}{(i,j)} dari Suatu Hasil Kali}

Pada bagian-bagian sebelumnya, kita menggunakan entri suatu matriks untuk menjelaskan cara kerja penjumlahan matriks dan perkalian skalar.
Kita juga dapat mempelajari perkalian matriks dengan menggunakan entri-entri matriks. 

Apa entri pada posisi $(i,j)$ dari $AB$? Entri tersebut berada pada baris $i$
dan kolom $j$ dari hasil kali $AB$.
 
Sekarang, jika $A$ berukuran $m \times n$ dan $B$ berukuran $n \times p$, kita mengetahui bahwa hasil kali $AB$ berbentuk
\begin{equation*}
\leftB
\begin{array}{cccc}
a_{11} & a_{12} & \cdots & a_{1n} \\
a_{21} & a_{22} & \cdots & a_{2n} \\
\vdots & \vdots &  \ddots & \vdots \\
a_{m1} & a_{m2} & \cdots & a_{mn}
\end{array}
\rightB \leftB
\begin{array}{cccccc}
b_{11} & b_{12} & \cdots & b_{1j} & \cdots & b_{1p} \\
b_{21} & b_{22} & \cdots & b_{2j} & \cdots & b_{2p} \\
\vdots & \vdots &  & \vdots & & \vdots\\
b_{n1} & b_{n2} & \cdots & b_{nj} & \cdots & b_{np}
\end{array}
\rightB 
\end{equation*}

Kolom $j$ dari $AB$ berbentuk
\begin{equation*}
\leftB
\begin{array}{cccc}
a_{11} & a_{12} & \cdots & a_{1n} \\
a_{21} & a_{22} & \cdots & a_{2n} \\
\vdots & \vdots & \ddots & \vdots \\
a_{m1} & a_{m2} & \cdots & a_{mn}
\end{array}
\rightB \leftB
\begin{array}{c}
b_{1j} \\
b_{2j} \\
\vdots \\
b_{nj}
\end{array}
\rightB
\end{equation*}
yang merupakan vektor kolom $m\times 1$. Vektor ini dihitung dengan
\begin{equation*}
b_{1j}
\leftB
\begin{array}{c}
a_{11} \\
a_{21} \\
\vdots \\
a_{m1}
\end{array}
\rightB + b_{2j}\leftB
\begin{array}{c}
a_{12} \\
a_{22} \\
\vdots \\
a_{m2}
\end{array}
\rightB +\cdots + b_{nj}\leftB
\begin{array}{c}
a_{1n} \\
a_{2n} \\
\vdots \\
a_{mn}
\end{array}
\rightB 
\end{equation*}

Oleh karena itu, entri pada posisi $(i,j)$ adalah entri pada baris $i$ dari vektor ini.
Entri tersebut dihitung dengan
\begin{equation*}
a_{i1}b_{1j}+a_{i2}b_{2j}+\cdots + a_{in}b_{nj}=\sum_{k=1}^{n}a_{ik}b_{kj}
\end{equation*}

Berikut adalah definisi formal entri pada posisi $(i,j)$ dari hasil kali matriks. 

\begin{definition}{Entri $(i,j)$ dari Suatu Hasil Kali}\label{def:ijentryofproduct}
Misalkan $A=\leftB a_{ij}\rightB $ adalah matriks $m\times n$ dan
$B=\leftB b_{ij}\rightB $ adalah matriks $n\times p$. Maka $AB$ adalah matriks
$m\times p$ dan entri-$\left( i, j \right)$ dari $AB$ \index{perkalian matriks!entri (i,j)} didefinisikan sebagai
\begin{equation*}
(AB)_{ij}=\sum_{k=1}^{n}a_{ik}b_{kj}  
\label{mat12}
\end{equation*}
Cara lain untuk menuliskannya adalah
\begin{equation*}
(AB)_{ij}=\leftB
\begin{array}{cccc}
a_{i1} & a_{i2} & \cdots & a_{in}
\end{array}
\rightB \leftB
\begin{array}{c}
b_{1j} \\
b_{2j} \\
\vdots \\
b_{nj}
\end{array}
\rightB
= 
a_{i1}b_{1j} + a_{i2}b_{2j} + \cdots + a_{in}b_{nj}
\end{equation*}
\end{definition}

Dengan kata lain, untuk mencari entri-$\left( i, j \right)$ dari hasil kali $AB$, atau $(AB)_{ij}$,
Anda mengalikan baris $i$ dari
$A,$ di sebelah kiri dengan kolom $j$ dari $B$. Untuk menyatakan $AB$ berdasarkan entri-entrinya, kita menulis $AB = \leftB (AB)_{ij} \rightB$.

Perhatikan contoh berikut. 

\begin{example}{Entri-Entri Suatu Hasil Kali}\label{exa:entriesofproduct1}
Hitung $AB$ jika memungkinkan. Jika dapat dilakukan, carilah entri-$\left( 3,2 \right)$ dari $AB$ menggunakan Definisi \ref{def:ijentryofproduct}.
\begin{equation*}
A = \leftB
\begin{array}{cc}
1 & 2 \\
3 & 1 \\
2 & 6
\end{array}
\rightB, B = \leftB
\begin{array}{ccc}
2 & 3 & 1 \\
7 & 6 & 2
\end{array}
\rightB
\end{equation*}
\end{example}

\begin{solution} Pertama, periksa apakah hasil kali dapat dilakukan. Bentuknya adalah $\left( 3\times
2\right) \left( 2\times 3\right) $, dan karena kedua bilangan di dalam sama, perkalian dapat dilakukan. Hasilnya seharusnya berupa matriks $3\times 3$.
Kita dapat terlebih dahulu menghitung $AB$:
\begin{equation*}
\leftB \leftB
\begin{array}{rr}
1 & 2 \\
3 & 1 \\
2 & 6
\end{array}
\rightB \leftB
\begin{array}{r}
2 \\
7
\end{array}
\rightB ,\leftB
\begin{array}{rr}
1 & 2 \\
3 & 1 \\
2 & 6
\end{array}
\rightB \leftB
\begin{array}{r}
3 \\
6
\end{array}
\rightB ,\leftB
\begin{array}{rr}
1 & 2 \\
3 & 1 \\
2 & 6
\end{array}
\rightB \leftB
\begin{array}{r}
1 \\
2
\end{array}
\rightB \rightB
\end{equation*}
dengan tanda koma memisahkan kolom-kolom pada hasil kali yang diperoleh. Jadi,
hasil kali di atas sama dengan
\begin{equation*}
 \leftB
\begin{array}{rrr}
16 & 15 & 5 \\
13 & 15 & 5 \\
46 & 42 & 14
\end{array}
\rightB 
\end{equation*}
yang merupakan matriks $3\times 3$, seperti yang diinginkan. Dengan demikian, entri-$\left( 3,2 \right)$ sama dengan 42.

Sekarang, dengan menggunakan Definisi
\ref{def:ijentryofproduct}, kita dapat memperoleh bahwa entri-$\left( 3,2 \right)$ sama dengan
\begin{eqnarray*}
\sum_{k=1}^{2}a_{3k}b_{k2} &=&a_{31}b_{12}+a_{32}b_{22} \\
&=&2\times 3+6\times 6=42\\
\end{eqnarray*}
Bila dibandingkan dengan hasil $AB$ di atas, nilai ini benar! 

Anda dapat menggunakan metode ini untuk memeriksa bahwa entri-entri $AB$ lainnya juga benar.
\end{solution}

Berikut contoh lainnya.

\begin{example}{Menentukan Entri-Entri Suatu Hasil Kali}\label{exa:entriesofproduct2}
Tentukan apakah hasil kali $AB$ terdefinisi. Jika ya, carilah entri-$\left( 2, 1 \right)$ dari hasil kali tersebut.
\begin{equation*}
A=
\leftB
\begin{array}{rrr}
2 & 3 & 1 \\
7 & 6 & 2 \\
0 & 0 & 0
\end{array}
\rightB,  B =\leftB
\begin{array}{rr}
1 & 2 \\
3 & 1 \\
2 & 6
\end{array}
\rightB 
\end{equation*}
\end{example}

\begin{solution} Hasil kali ini berbentuk $\left( 3\times
3\right) \left( 3\times 2\right) $. Kedua bilangan di tengah sama, sehingga
matriks-matriks tersebut konformabel dan hasil kalinya dapat dihitung.  

Kita ingin mencari entri-$\left( 2, 1 \right)$ dari $AB$, yaitu entri pada baris kedua dan kolom pertama hasil kali.
Kita akan menggunakan Definisi \ref{def:ijentryofproduct}, yang menyatakan
\begin{equation*}
(AB)_{ij}=\sum_{k=1}^{n}a_{ik}b_{kj}
\end{equation*}
Dalam hal ini, $n=3$, $i=2$, dan $j=1$. Jadi, entri-$\left( 2, 1 \right)$ diperoleh dengan menghitung
\begin{equation*}
 (AB)_{21} = \sum_{k=1}^{3}a_{2k}b_{k1} = 
 \leftB
\begin{array}{ccc}
a_{21} & a_{22} & a_{23}
\end{array}
\rightB \leftB
\begin{array}{c}
b_{11} \\
b_{21} \\
b_{31}
\end{array}
\rightB  
\end{equation*}
Dengan menyubstitusikan nilai-nilai yang sesuai, hasil kali ini menjadi
\begin{equation*}
\leftB
\begin{array}{ccc}
a_{21} & a_{22} & a_{23}
\end{array}
\rightB \leftB
\begin{array}{c}
b_{11} \\
b_{21} \\
b_{31}
\end{array}
\rightB 
=
\leftB
\begin{array}{ccc}
7 & 6 & 2
\end{array}
\rightB \leftB
\begin{array}{c}
1 \\
3 \\
2
\end{array}
\rightB
=
1 \times 7 + 3 \times 6 + 2 \times 2
=
29
\end{equation*}

Jadi, $(AB)_{21} = 29$.

Luangkan waktu sejenak untuk mencari beberapa entri $AB$ lainnya. Anda dapat mengalikan matriks-matriks tersebut untuk memeriksa bahwa jawaban Anda benar.
Hasil kali $AB$ diberikan oleh
\begin{equation*}
AB = \leftB
\begin{array}{cc}
13 & 13 \\
29 & 32 \\
0 & 0
\end{array}
\rightB 
\end{equation*}
\end{solution}
